%قبل از شروع مسئله ترجیحاً چیزی ننویسید. اما اگر خواستید بنویسید، در قالب لیستی (فقرات و شمارش) بنویسید.
%برای اضافه کردن نمره، بعدی از شروع‌مسئله، [x] را اضافه کنید و به جای x  نمره مورد نظر را قرار دهید. برای زیرمسئله‌ها هم به همین صورت است.
%\شروع‌مسئله[10]
\شروع‌مسئله [40]
%\متن‌سیاه{سوال عملی اول:}\\
%\متن‌سیاه{دراپ‌تست}\\

جیسون که یکی از مأموران اطلاعاتی سازمان سیا است به مأموریتی در کشوری آسیایی فرستاده شده است. هدف او از این مأموریت پیداکردن طعم آدامس مورد علاقه رئیس جمهور است. 
جیسون پس از انجام این مأموریت قصد دارد تا این طعم را به صورت مخفیانه برای سازمان سیا ارسال کند. اما چون رمزنگار بسیار ضعیفی است تصمیم می‌گیرد که این طعم را در دو رشته تصادفی از حروف انگلیسی جاگذاری کند (طعم آدامس را به صورت پیوسته در نقطه‌ای تصادفی در دو رشته قرار می‌هد.)  و تضمین می‌کند که دو رشته هیچ زیررشته مشترکی با طول بیشتر مساوی از طول طعم  آدامس  ندارند و آن دو رشته را برای سازمان سیا با ایمیل ارسال کند. خوشبختانه این راه رمزنگاری او لو می‌رود و سعید قصد دارد این طعم را بیابد. سعید موفق شده است که ایمیل جیسون را هک کند. سعید نیاز به برنامه‌ای دارد با دادن دو رشته طعم آدامس رئیس جهمور را بدست بیاورد.
به سعید کمک کنید تا این طعم را بیابد.

\متن‌سیاه{ورودی:}\\
دو رشته که همان دو رشته‌ای هستند که جیسون برای سازمان سیا ارسال کرده است.
طول دو رشته حداکثر
$5 \times 10^3$
است.

\متن‌سیاه{خروجی:}\\
طعم آدامس مورد علاقه رئیس جمهور که در دو رشته جاگذاری شده 
است.

\متن‌سیاه{محدودیت‌ها:}\\
رشته‌ها از حروف کوچک و بزرگ انگلیسی تشکیل شده‌اند.
استفاده از داده ساختار dictionary پایتون مجاز نمی باشد. برای مقایسه رشته ها اجازه استفاده از عملگر == را ندارید. همچنین اجازه استفاده از داده ساختارهایی که در کلاس درس تدریس نشده اند را ندارید.

\متن‌سیاه{ورودی و خروجی نمونه:}\\

تست اول
\begin{latin}
	\begin{tabularx}{\textwidth}{|X | X|}
		\hline
		\textbf{stdin} &
		\textbf{stdout}\\
		\hline
		ssdcmsdersddlimooodmqccsernmf & limoo \\
		dsflimoosmcneerbfdser & \\
		\hline
	\end{tabularx}
\end{latin}


\solution{


اگر بیشینه‌‌ی طول رشته‌ها را $n$ بگیریم، ابتدا برای هر یک از رشته‌ها به کمک الگوریتم rolling hash  در $O(n)$ جدول هش مربوط به $i$ کارارکتر اول رشته را تشکیل می‌دهیم. پس از ساخت این آرایه می‌توانیم که مقدار هش هر زیربازه‌ای از این رشته را در $O(1)$ بیابیم.

حال روی طول بزرگترین زیررشته‌ی مشترک جست‌وجوی دودویی می‌کنیم، با در نظر گرفتن طول $l$ برای بررسی در جست‌وجو، همه‌ی $l$ در رشته‌ی اول را در $O(n-l)$ در یک جدول درهم‌سازی ریخته، سپس همه‌ی زیر رشته‌های به طول $l$ از رشته‌ی دوم را در این جدول جست‌وجو می‌کنیم. از آنجایی که هر جست‌وجو از $O(1)$ است، روی‌ هم $O(n-l)$ هزینه کرده‌ایم.
حال اگر یک از زیررشته‌های رشته‌ی دوم در جدول موجود بود به این معنی است که طول بزرگترین زیررشته بیشتر از $l$ است، $l$ را افزایش می‌دهیم در غیر این صورت کاهش می‌دهیم.



‌\\کد نمونه (با تشکر از ارشیا سلطانی) :

\definecolor{codegreen}{rgb}{0,0.6,0}
\definecolor{codegray}{rgb}{0.5,0.5,0.5}
\definecolor{codepurple}{rgb}{0.58,0,0.82}
\definecolor{backcolour}{rgb}{0.95,0.95,0.92}

\lstdefinestyle{mystyle}{
    backgroundcolor=\color{backcolour},   
    commentstyle=\color{codegreen},
    keywordstyle=\color{magenta},
    numberstyle=\tiny\color{codegray},
    stringstyle=\color{codepurple},
    basicstyle=\ttfamily\footnotesize,
    breakatwhitespace=false,         
    breaklines=true,                 
    captionpos=b,                    
    keepspaces=true,                 
    numbers=left,                    
    numbersep=5pt,                  
    showspaces=false,                
    showstringspaces=false,
    showtabs=false,                  
    tabsize=2
}

\lstset{style=mystyle}

\lr{\lstinputlisting[language=python]{src/adams.py}}


}

%}
% صورت اصلی مسئله را در این‌جا بنویسید.
%اگر مسئله چند قسمت داشت به این صورت عمل کنید:
%\شروع{زیرمسائل}
%\زیرمسئله[0] % مسئلۀ ۱آ
%\زیرمسئله % مسئلۀ ۱ب
%\زیرمسئله % مسئلۀ ۱پ
%\زیرمسئله % مسئلۀ ۱ت
%\پایان{زیرمسائل}



%%%%%%%%%%%%%%%%%%%%%%%%%%%%%%%%%%%%%%%%%%%%%%%%%%%%%
% برخی دستورات مفید برای تمرین‌ها. %
%%%%%%%%%%%%%%%%%%%%%%%%%%%%%%%%%%%%%%%%%%%%%%%%%%%%%

% فرمول‌های نمونه:
%$f(x) = \Theta(x)$
%$T(x, y) \leq \log(x) + 2^y + \binom{2n}{n}$
% {\tt code\_function}


% لیست موارد بدون شماره‌گذاری:
%\شروع{فقرات}
%    \فقره مورد اول 
%        \شروع{فقرات}
%            \فقره مورد اول داخلی 
%                \شروع{فقرات}
%                    \فقره اولین مورد خیلی داخلی 
%                    \فقره دومین مورد خیلی داخلی 
%                \پایان{فقرات}
%            \فقره مورد دوم داخلی 
%        \پایان{فقرات}
%    \فقره مورد دوم 
%\پایان{فقرات}

% لیست موارد شماره‌گذاری شده:
%\شروع{شمارش}
%    \فقره یک مورد 
%    \فقره یک مورد دیگر 
%    \فقره باز هم یک مورد دیگر
%\پایان{شمارش}

% برای نوشتن یک فرمول چند خطی که یک شماره بخوره. برای حالت‌های دیگر به سایت زیر مراجعه کنید:
% https://www.overleaf.com/learn/latex/Aligning_equations_with_amsmath
%\begin{align} 
%    \begin{split}
%        (x+y)^3 &= (x+y)^2(x+y) \\
%                &= (x^2+2xy+y^2)(x+y) \\
%                &= (x^3+2x^2y+xy^2) + (x^2y+2xy^2+y^3) \\
%                &= x^3+3x^2y+3xy^2+y^3
%    \end{split}                                 
%\end{align}

% نوشتن ماتریس و آرایه دوبعدی
%\begin{align}
%    A = 
%    \begin{bmatrix}
%        A_{11} & A_{21} \\
%        A_{21} & A_{22}
%    \end{bmatrix}
%\end{align}

% اضافه کردن تصویر با شرح. برای اضافه کردن تصویر بدون شرح، فقط دستور درج‌تصویر را استفاده کنید.
%\شروع{شکل}[h!]
%\تنظیم‌ازوسط
%  \درج‌تصویر[پهنا=0.5\پهنای‌متن]{img.jpg}
%\شرح{شرح تصویر}
%\برچسب{تصویر:۱}
%\پایان{شکل}
%\پایان{زیرمسائل}

%کد زیر را ببینید.
%%\lstinputlisting{source_filename.py}
%\begin{latin}
%\begin{lstlisting}
%# Import the modules
%import sys
%import random
%
%ans = True
%
%while ans:
%question = raw_input("Ask the magic 8 ball a question: (press enter to quit) ")
%
%answers = random.randint(1,8)
%
%if question == "":
%sys.exit()
%
%elif answers == 1:
%print "It is certain"
%
%elif answers == 2:
%print "Outlook good"
%
%elif answers == 3:
%print "You may rely on it"
%
%elif answers == 4:
%print "Ask again later"
%
%elif answers == 5:
%print "Concentrate and ask again"
%
%elif answers == 6:
%print "Reply hazy, try again"
%
%elif answers == 7:
%print "My reply is no"
%
%elif answers == 8:
%print "My sources say no"
%\end{lstlisting}
%\end{latin}
